% ======================================================================
% Navegacao
%   \bloco[curto]{Titulo}          bloco do deck (vira \part)
%   \modulo[curto]{rotulo}{Titulo} modulo (vira \section + menu do modulo)
%   \emconstrucao[curto]{rotulo}{Titulo}
%                                  destino ainda nao escrito
%   Barra lateral: bloco atual, seus modulos e, expandido, o modulo atual
% ======================================================================

% Rotulo do modulo corrente (usado pelo botao "Modulo" da barra lateral)
\providecommand{\moduloatual}{menu}

% ---------- Botoes ----------
\newcommand{\irpara}[2]{\hyperlink{#1}{\beamergotobutton{#2}}}
\newcommand{\voltar}{\Acrobatmenu{GoBack}{\beamerreturnbutton{Voltar}}}
\newcommand{\verrevisao}[1]{%
  \hyperlink{#1}{\beamergotobutton{Revisão \MakeUppercase{#1}}}}

% ---------- Estrutura ----------
\NewDocumentCommand{\bloco}{O{#2} m}{%
  \part[#1]{#2}}

\NewDocumentCommand{\modulo}{O{#3} m m}{%
  \gdef\moduloatual{#2}%
  \section[#1]{#3}%
  \begin{frame}[label=#2]{#3}
    \begin{multicols}{2}
      \small
      \tableofcontents[currentsection, hideothersubsections, sectionstyle=hide/hide]
    \end{multicols}
  \end{frame}}

\NewDocumentCommand{\emconstrucao}{O{#3} m m}{%
  \gdef\moduloatual{#2}%
  \section[#1]{#3}%
  \begin{frame}[label=#2]{#3}
    \centering
    \vfill
    {\Large\color{labsemGrafite!60} Em construção}
    \vfill
    \irpara{menu}{Menu}
  \end{frame}}

% ---------- Botao grande dos menus: \botaomenu{rotulo}{Texto} ----------
\newcommand{\botaomenu}[2]{%
  \hyperlink{#1}{%
    \tikz[baseline=(b.base)]
      \node[draw=labsemMarinho, fill=labsemGelo, rounded corners=3pt,
            minimum width=0.46\textwidth, minimum height=1.7em,
            text=labsemMarinho, font=\small\bfseries, align=center] (b) {#2};}}
